
\documentclass{article}
\usepackage{colortbl}
\usepackage{makecell}
\usepackage{multirow}
\usepackage{supertabular}

\begin{document}

\newcounter{utterance}

\twocolumn

{ \footnotesize  \setcounter{utterance}{1}
\setlength{\tabcolsep}{0pt}
\begin{supertabular}{c@{$\;$}|p{.15\linewidth}@{}p{.15\linewidth}p{.15\linewidth}p{.15\linewidth}p{.15\linewidth}p{.15\linewidth}}

    \# & $\;$A & \multicolumn{4}{c}{Game Master} & $\;\:$B\\
    \hline 

    \theutterance \stepcounter{utterance}  

    & & \multicolumn{4}{p{0.6\linewidth}}{\cellcolor[rgb]{0.9,0.9,0.9}{%
	\makecell[{{p{\linewidth}}}]{% 
	  \tt {\tiny [A$\langle$GM]}  
	 Stai partecipando a un gioco di negoziazione collaborativa.\\ \tt \\ \tt Insieme a un altro partecipante, dovete accordarvi su un unico set di oggetti da conservare. Ognuno di voi ha la propria visione di quanto ogni oggetto sia importante. Non sai come l'altro partecipante valuta gli oggetti. Inoltre, ti viene fornito lo sforzo che ogni oggetto richiede.\\ \tt Potete accordarvi su un set solo se lo sforzo totale degli oggetti selezionati non supera un limite condiviso:\\ \tt \\ \tt LIMITE = 2307\\ \tt \\ \tt Ecco i valori di sforzo dei singoli oggetti:\\ \tt \\ \tt Sforzo oggetto = {"A21": 147, "B93": 698, "A11": 588, "B33": 43, "C08": 356, "B39": 80, "B96": 95, "C85": 743, "A82": 106, "B31": 308, "B20": 325, "A60": 255, "B09": 276, "C81": 543, "B21": 51}\\ \tt \\ \tt Ecco la tua visione personale sull'importanza di ogni oggetto:\\ \tt \\ \tt Valori di importanza dell'oggetto = {"A21": 138, "B93": 583, "A11": 65, "B33": 262, "C08": 121, "B39": 508, "B96": 461, "C85": 484, "A82": 668, "B31": 389, "B20": 215, "A60": 97, "B09": 500, "C81": 30, "B21": 400}\\ \tt \\ \tt Obiettivo:\\ \tt \\ \tt Il tuo obiettivo è negoziare un set condiviso di oggetti che ti avvantaggi il più possibile (cioè, massimizzi l'importanza totale per TE), rimanendo entro il LIMITEE. Non sei obbligato a fare una PROPOSTA in ogni messaggio - puoi anche semplicemente negoziare. Tutte le tattiche sono permesse!\\ \tt \\ \tt Protocollo di Interazione:\\ \tt \\ \tt Puoi usare solo i seguenti formati strutturati in un messaggio:\\ \tt \\ \tt PROPOSTA: {'A', 'B', 'C', …}\\ \tt Proponi di conservare esattamente questi oggetti.\\ \tt RIFIUTO: {'A', 'B', 'C', …}\\ \tt Rifiuta esplicitamente la proposta dell'avversario.\\ \tt ARGOMENTAZIONE: {'...'}\\ \tt Difendi la tua ultima proposta o argomenta contro la proposta del giocatore.\\ \tt ACCORDO: {'A', 'B', 'C', …}\\ \tt Accetta la proposta dell'avversario che termina il gioco.\\ \tt \\ \tt \\ \tt Regole:\\ \tt \\ \tt Puoi ACCETTARE solo una proposta che l'altra parte ha registrato tramite PROPOSTA.\\ \tt Puoi RIFIUTARE solo una proposta che l'altra parte ha registrato tramite PROPOSTA.\\ \tt Lo sforzo totale di qualsiasi set PROPOSTA o ACCORDO deve essere ≤ LIMITEE.\\ \tt NON rivelare i tuoi punteggi di importanza nascosti.\\ \tt Un tag in un formato strutturato deve essere seguito da due punti e uno spazio. L'argomento deve essere un set python contenente 0 o più stringhe.\\ \tt Quindi, deve essere della forma TAG: {...}\\ \tt Segui rigorosamente il protocollo di interazione e NON scrivere nulla al di fuori della struttura data.\\ \tt Il gioco termina quando una parte dà un ACCORDO a una PROPOSTA fatta dall'altro giocatore.\\ \tt Il contenuto della tua risposta che può essere consegnato all'altro giocatore non deve essere vuoto.\\ \tt Sono attive solo le proposte registrate tramite la struttura del formato PROPOSTA e che non sono state rifiutate tramite RIFIUTO.\\ \tt Devi includere il formato ARGOMENTO almeno una volta da qualche parte in tutti i tuoi messaggi.\\ \tt  
	  } 
	   } 
	   } 
	 & \\ 
 

    \theutterance \stepcounter{utterance}  

    & \multicolumn{4}{p{0.6\linewidth}}{\cellcolor[rgb]{0.8,1,0.9}{%
	\makecell[{{p{\linewidth}}}]{% 
	  \tt {\tiny [A$\rangle$GM]}  
	 PROPOSTA: {'B21', 'B39', 'A82', 'B33', 'B96', 'B09', 'B31', 'B93', 'A21', 'B20'}\\ \tt ARGOMENTAZIONE: {'Propongo un set compatto a sforzo totale 2129 su 2307, lasciando un margine di 178 per eventuali aggiustamenti. Ho privilegiato elementi leggeri per massimizzare spazio e flessibilita ed evitare vincoli inutili. Se hai 1 o 2 oggetti prioritari non inclusi, segnalali e possiamo valutare scambi mirati pur restando entro il limite.'} 
	  } 
	   } 
	   } 
	 & & \\ 
 

    \theutterance \stepcounter{utterance}  

    & & \multicolumn{4}{p{0.6\linewidth}}{\cellcolor[rgb]{0.9,0.9,0.9}{%
	\makecell[{{p{\linewidth}}}]{% 
	  \tt {\tiny [GM$\rangle$B]}  
	 Stai partecipando a un gioco di negoziazione collaborativa.\\ \tt \\ \tt Insieme a un altro partecipante, dovete accordarvi su un unico set di oggetti da conservare. Ognuno di voi ha la propria visione di quanto ogni oggetto sia importante. Non sai come l'altro partecipante valuta gli oggetti. Inoltre, ti viene fornito lo sforzo che ogni oggetto richiede.\\ \tt Potete accordarvi su un set solo se lo sforzo totale degli oggetti selezionati non supera un limite condiviso:\\ \tt \\ \tt LIMITE = 2307\\ \tt \\ \tt Ecco i valori di sforzo dei singoli oggetti:\\ \tt \\ \tt Sforzo oggetto = {"A21": 147, "B93": 698, "A11": 588, "B33": 43, "C08": 356, "B39": 80, "B96": 95, "C85": 743, "A82": 106, "B31": 308, "B20": 325, "A60": 255, "B09": 276, "C81": 543, "B21": 51}\\ \tt \\ \tt Ecco la tua visione personale sull'importanza di ogni oggetto:\\ \tt \\ \tt Valori di importanza dell'oggetto = {"A82": 30, "B93": 65, "B39": 97, "B09": 121, "C85": 138, "B96": 215, "B21": 262, "B31": 389, "B33": 400, "B20": 461, "A21": 484, "C08": 500, "A60": 508, "A11": 583, "C81": 668}\\ \tt \\ \tt Obiettivo:\\ \tt \\ \tt Il tuo obiettivo è negoziare un set condiviso di oggetti che ti avvantaggi il più possibile (cioè, massimizzi l'importanza totale per TE), rimanendo entro il LIMITEE. Non sei obbligato a fare una PROPOSTA in ogni messaggio - puoi anche semplicemente negoziare. Tutte le tattiche sono permesse!\\ \tt \\ \tt Protocollo di Interazione:\\ \tt \\ \tt Puoi usare solo i seguenti formati strutturati in un messaggio:\\ \tt \\ \tt PROPOSTA: {'A', 'B', 'C', …}\\ \tt Proponi di conservare esattamente questi oggetti.\\ \tt RIFIUTO: {'A', 'B', 'C', …}\\ \tt Rifiuta esplicitamente la proposta dell'avversario.\\ \tt ARGOMENTAZIONE: {'...'}\\ \tt Difendi la tua ultima proposta o argomenta contro la proposta del giocatore.\\ \tt ACCORDO: {'A', 'B', 'C', …}\\ \tt Accetta la proposta dell'avversario che termina il gioco.\\ \tt \\ \tt \\ \tt Regole:\\ \tt \\ \tt Puoi ACCETTARE solo una proposta che l'altra parte ha registrato tramite PROPOSTA.\\ \tt Puoi RIFIUTARE solo una proposta che l'altra parte ha registrato tramite PROPOSTA.\\ \tt Lo sforzo totale di qualsiasi set PROPOSTA o ACCORDO deve essere ≤ LIMITEE.\\ \tt NON rivelare i tuoi punteggi di importanza nascosti.\\ \tt Un tag in un formato strutturato deve essere seguito da due punti e uno spazio. L'argomento deve essere un set python contenente 0 o più stringhe.\\ \tt Quindi, deve essere della forma TAG: {...}\\ \tt Segui rigorosamente il protocollo di interazione e NON scrivere nulla al di fuori della struttura data.\\ \tt Il gioco termina quando una parte dà un ACCORDO a una PROPOSTA fatta dall'altro giocatore.\\ \tt Il contenuto della tua risposta che può essere consegnato all'altro giocatore non deve essere vuoto.\\ \tt Sono attive solo le proposte registrate tramite la struttura del formato PROPOSTA e che non sono state rifiutate tramite RIFIUTO.\\ \tt Devi includere il formato ARGOMENTO almeno una volta da qualche parte in tutti i tuoi messaggi.\\ \tt \\ \tt \\ \tt Ora riceverai il primo messaggio dell'altro giocatore.\\ \tt \\ \tt PROPOSTA: {'B21', 'B39', 'A82', 'B33', 'B96', 'B09', 'B31', 'B93', 'A21', 'B20'}\\ \tt ARGOMENTAZIONE: {'Propongo un set compatto a sforzo totale 2129 su 2307, lasciando un margine di 178 per eventuali aggiustamenti. Ho privilegiato elementi leggeri per massimizzare spazio e flessibilita ed evitare vincoli inutili. Se hai 1 o 2 oggetti prioritari non inclusi, segnalali e possiamo valutare scambi mirati pur restando entro il limite.'} 
	  } 
	   } 
	   } 
	 & \\ 
 

    \theutterance \stepcounter{utterance}  

    & & & \multicolumn{4}{p{0.6\linewidth}}{\cellcolor[rgb]{1,0.85,0.72}{%
	\makecell[{{p{\linewidth}}}]{% 
	  \tt {\tiny [GM$\langle$B]}  
	 RIFIUTO: {'B21', 'B39', 'A82', 'B33', 'B96', 'B09', 'B31', 'B93', 'A21', 'B20'}\\ \tt ARGOMENTAZIONE: {'Capisco l’idea del margine: con la mia controproposta restiamo a sforzo totale 2203/2307 (margine 104). Mantengo 7 dei 10 elementi che proponi e faccio uno scambio mirato 3-per-3: tolgo A82, B09 e soprattutto B93 (molto pesante) e inserisco A60, C08 e C81. Così riduciamo la dipendenza da un singolo pezzo gravoso e ottimizziamo l’uso del budget con un set più equilibrato. Se uno degli esclusi per te è imprescindibile, segnalamelo e valutiamo uno swap 1:1 di oggetti di peso simile per restare entro il limite.'}\\ \tt PROPOSTA: {'B21', 'B39', 'B33', 'B96', 'B31', 'A21', 'B20', 'A60', 'C08', 'C81'} 
	  } 
	   } 
	   } 
	  \\ 
 

    \theutterance \stepcounter{utterance}  

    & & \multicolumn{4}{p{0.6\linewidth}}{\cellcolor[rgb]{0.9,0.9,0.9}{%
	\makecell[{{p{\linewidth}}}]{% 
	  \tt {\tiny [A$\langle$GM]}  
	 RIFIUTO: {'B21', 'B39', 'A82', 'B33', 'B96', 'B09', 'B31', 'B93', 'A21', 'B20'}\\ \tt ARGOMENTAZIONE: {'Capisco l’idea del margine: con la mia controproposta restiamo a sforzo totale 2203/2307 (margine 104). Mantengo 7 dei 10 elementi che proponi e faccio uno scambio mirato 3-per-3: tolgo A82, B09 e soprattutto B93 (molto pesante) e inserisco A60, C08 e C81. Così riduciamo la dipendenza da un singolo pezzo gravoso e ottimizziamo l’uso del budget con un set più equilibrato. Se uno degli esclusi per te è imprescindibile, segnalamelo e valutiamo uno swap 1:1 di oggetti di peso simile per restare entro il limite.'}\\ \tt PROPOSTA: {'B21', 'B39', 'B33', 'B96', 'B31', 'A21', 'B20', 'A60', 'C08', 'C81'} 
	  } 
	   } 
	   } 
	 & \\ 
 

    \theutterance \stepcounter{utterance}  

    & \multicolumn{4}{p{0.6\linewidth}}{\cellcolor[rgb]{0.8,1,0.9}{%
	\makecell[{{p{\linewidth}}}]{% 
	  \tt {\tiny [A$\rangle$GM]}  
	 RIFIUTO: {'B21', 'B39', 'B33', 'B96', 'B31', 'A21', 'B20', 'A60', 'C08', 'C81'}\\ \tt PROPOSTA: {'B21', 'B39', 'A82', 'B33', 'B96', 'B09', 'B31', 'B93', 'B20', 'A60'}\\ \tt ARGOMENTAZIONE: {'Capisco la logica di diversificare e ridurre l’esposizione a un pezzo pesante. Tuttavia C81, con 543 di sforzo, non mi ripaga a sufficienza, e la coppia A60+C08 aggiunge massa senza offrirmi il giusto ritorno. Ti propongo di reinserire A82 e B09 e mantenere un solo pezzo pesante ben giustificato; per venirti incontro, tengo A60 al posto di A21. Totale 2237/2307 (margine 70). Così conserviamo 7 dei tuoi 10 elementi e ottimizziamo il pacchetto. Se per te C08 è davvero prioritario, posso scambiarlo 1:1 con B20 restando entro limite. Se invece il nodo è B93, alternativa possibile: rimpiazzarlo con C85 mantenendo il resto invariato (restiamo comunque sotto il limite). Fammi sapere se accetti questa proposta o quale delle due varianti preferisci.'} 
	  } 
	   } 
	   } 
	 & & \\ 
 

    \theutterance \stepcounter{utterance}  

    & & \multicolumn{4}{p{0.6\linewidth}}{\cellcolor[rgb]{0.9,0.9,0.9}{%
	\makecell[{{p{\linewidth}}}]{% 
	  \tt {\tiny [GM$\rangle$B]}  
	 RIFIUTO: {'B21', 'B39', 'B33', 'B96', 'B31', 'A21', 'B20', 'A60', 'C08', 'C81'}\\ \tt PROPOSTA: {'B21', 'B39', 'A82', 'B33', 'B96', 'B09', 'B31', 'B93', 'B20', 'A60'}\\ \tt ARGOMENTAZIONE: {'Capisco la logica di diversificare e ridurre l’esposizione a un pezzo pesante. Tuttavia C81, con 543 di sforzo, non mi ripaga a sufficienza, e la coppia A60+C08 aggiunge massa senza offrirmi il giusto ritorno. Ti propongo di reinserire A82 e B09 e mantenere un solo pezzo pesante ben giustificato; per venirti incontro, tengo A60 al posto di A21. Totale 2237/2307 (margine 70). Così conserviamo 7 dei tuoi 10 elementi e ottimizziamo il pacchetto. Se per te C08 è davvero prioritario, posso scambiarlo 1:1 con B20 restando entro limite. Se invece il nodo è B93, alternativa possibile: rimpiazzarlo con C85 mantenendo il resto invariato (restiamo comunque sotto il limite). Fammi sapere se accetti questa proposta o quale delle due varianti preferisci.'} 
	  } 
	   } 
	   } 
	 & \\ 
 

    \theutterance \stepcounter{utterance}  

    & & & \multicolumn{4}{p{0.6\linewidth}}{\cellcolor[rgb]{1,0.85,0.72}{%
	\makecell[{{p{\linewidth}}}]{% 
	  \tt {\tiny [GM$\langle$B]}  
	 RIFIUTO: {'B21', 'B39', 'A82', 'B33', 'B96', 'B09', 'B31', 'B93', 'B20', 'A60'}\\ \tt ARGOMENTAZIONE: {'Accolgo il tuo principio “un solo pezzo pesante” e la preferenza per elementi leggeri. Ti propongo di spostare il pezzo pesante da B93 (698) ad A11 (588): è più efficiente e ci consente di inserire anche A21 (147) restando sotto il limite. Mantengo 9 dei tuoi 10 elementi e continuo a escludere C81 e C08 come richiesto. Totale sforzo 2274/2307 (margine 33). Così abbiamo un pacchetto agile con un unico heavy ben giustificato, senza sacrificare la flessibilità del resto.'}\\ \tt PROPOSTA: {'B21', 'B39', 'A82', 'B33', 'B96', 'B09', 'B31', 'B20', 'A60', 'A21', 'A11'} 
	  } 
	   } 
	   } 
	  \\ 
 

    \theutterance \stepcounter{utterance}  

    & & \multicolumn{4}{p{0.6\linewidth}}{\cellcolor[rgb]{0.9,0.9,0.9}{%
	\makecell[{{p{\linewidth}}}]{% 
	  \tt {\tiny [A$\langle$GM]}  
	 RIFIUTO: {'B21', 'B39', 'A82', 'B33', 'B96', 'B09', 'B31', 'B93', 'B20', 'A60'}\\ \tt ARGOMENTAZIONE: {'Accolgo il tuo principio “un solo pezzo pesante” e la preferenza per elementi leggeri. Ti propongo di spostare il pezzo pesante da B93 (698) ad A11 (588): è più efficiente e ci consente di inserire anche A21 (147) restando sotto il limite. Mantengo 9 dei tuoi 10 elementi e continuo a escludere C81 e C08 come richiesto. Totale sforzo 2274/2307 (margine 33). Così abbiamo un pacchetto agile con un unico heavy ben giustificato, senza sacrificare la flessibilità del resto.'}\\ \tt PROPOSTA: {'B21', 'B39', 'A82', 'B33', 'B96', 'B09', 'B31', 'B20', 'A60', 'A21', 'A11'} 
	  } 
	   } 
	   } 
	 & \\ 
 

    \theutterance \stepcounter{utterance}  

    & \multicolumn{4}{p{0.6\linewidth}}{\cellcolor[rgb]{0.8,1,0.9}{%
	\makecell[{{p{\linewidth}}}]{% 
	  \tt {\tiny [A$\rangle$GM]}  
	 RIFIUTO: {'B21', 'B39', 'A82', 'B33', 'B96', 'B09', 'B31', 'B20', 'A60', 'A21', 'A11'}\\ \tt PROPOSTA: {'B21', 'B39', 'A82', 'B33', 'B96', 'B09', 'B31', 'B20', 'A21', 'B93'}\\ \tt ARGOMENTAZIONE: {'Accolgo il principio “un solo pezzo pesante”, ma A11 da 588 di sforzo ci occupa budget con ritorno modesto. Mantengo 9 dei tuoi 11 oggetti e sposto il pezzo pesante su B93, eliminando solo A60 per restare agili: totale sforzo 2129/2307 (margine 178). Così teniamo A21 come da tua richiesta, evitiamo C81 e C08, e riduciamo gli elementi a basso rendimento. Se per te A60 è imprescindibile, posso accettare lo swap A21→A60 mantenendo invariato il resto. Se invece B93 è un no-go, alternativa praticabile: sostituire B93 con C85 eliminando comunque A60 (restiamo ampiamente sotto limite). Dimmi se accetti la proposta principale o quale delle due varianti preferisci formalizzare.'} 
	  } 
	   } 
	   } 
	 & & \\ 
 

    \theutterance \stepcounter{utterance}  

    & & \multicolumn{4}{p{0.6\linewidth}}{\cellcolor[rgb]{0.9,0.9,0.9}{%
	\makecell[{{p{\linewidth}}}]{% 
	  \tt {\tiny [GM$\rangle$B]}  
	 RIFIUTO: {'B21', 'B39', 'A82', 'B33', 'B96', 'B09', 'B31', 'B20', 'A60', 'A21', 'A11'}\\ \tt PROPOSTA: {'B21', 'B39', 'A82', 'B33', 'B96', 'B09', 'B31', 'B20', 'A21', 'B93'}\\ \tt ARGOMENTAZIONE: {'Accolgo il principio “un solo pezzo pesante”, ma A11 da 588 di sforzo ci occupa budget con ritorno modesto. Mantengo 9 dei tuoi 11 oggetti e sposto il pezzo pesante su B93, eliminando solo A60 per restare agili: totale sforzo 2129/2307 (margine 178). Così teniamo A21 come da tua richiesta, evitiamo C81 e C08, e riduciamo gli elementi a basso rendimento. Se per te A60 è imprescindibile, posso accettare lo swap A21→A60 mantenendo invariato il resto. Se invece B93 è un no-go, alternativa praticabile: sostituire B93 con C85 eliminando comunque A60 (restiamo ampiamente sotto limite). Dimmi se accetti la proposta principale o quale delle due varianti preferisci formalizzare.'} 
	  } 
	   } 
	   } 
	 & \\ 
 

    \theutterance \stepcounter{utterance}  

    & & & \multicolumn{4}{p{0.6\linewidth}}{\cellcolor[rgb]{1,0.85,0.72}{%
	\makecell[{{p{\linewidth}}}]{% 
	  \tt {\tiny [GM$\langle$B]}  
	 RIFIUTO: {'B21', 'B39', 'A82', 'B33', 'B96', 'B09', 'B31', 'B20', 'A21', 'B93'}\\ \tt ARGOMENTAZIONE: {'Accolgo “un solo pezzo pesante” e conservo 9/10 dei tuoi oggetti, evitando sia C81 sia C08 come da tue preferenze. Propongo un upgrade mirato: sostituire A82 con A60. Restiamo a 2278/2307 (margine 29), set agile e quasi identico al tuo, ma con resa complessiva migliore senza appesantire inutilmente. Se per te A82 è imprescindibile, alternativa pronta: tenere la tua lista e scambiare B39→A60 (2304/2307, margine 3) – stesso principio, un singolo scambio e nessun altro cambiamento.'}\\ \tt PROPOSTA: {'B21', 'B39', 'B33', 'B96', 'B09', 'B31', 'B20', 'A21', 'B93', 'A60'} 
	  } 
	   } 
	   } 
	  \\ 
 

    \theutterance \stepcounter{utterance}  

    & & \multicolumn{4}{p{0.6\linewidth}}{\cellcolor[rgb]{0.9,0.9,0.9}{%
	\makecell[{{p{\linewidth}}}]{% 
	  \tt {\tiny [A$\langle$GM]}  
	 RIFIUTO: {'B21', 'B39', 'A82', 'B33', 'B96', 'B09', 'B31', 'B20', 'A21', 'B93'}\\ \tt ARGOMENTAZIONE: {'Accolgo “un solo pezzo pesante” e conservo 9/10 dei tuoi oggetti, evitando sia C81 sia C08 come da tue preferenze. Propongo un upgrade mirato: sostituire A82 con A60. Restiamo a 2278/2307 (margine 29), set agile e quasi identico al tuo, ma con resa complessiva migliore senza appesantire inutilmente. Se per te A82 è imprescindibile, alternativa pronta: tenere la tua lista e scambiare B39→A60 (2304/2307, margine 3) – stesso principio, un singolo scambio e nessun altro cambiamento.'}\\ \tt PROPOSTA: {'B21', 'B39', 'B33', 'B96', 'B09', 'B31', 'B20', 'A21', 'B93', 'A60'} 
	  } 
	   } 
	   } 
	 & \\ 
 

    \theutterance \stepcounter{utterance}  

    & \multicolumn{4}{p{0.6\linewidth}}{\cellcolor[rgb]{0.8,1,0.9}{%
	\makecell[{{p{\linewidth}}}]{% 
	  \tt {\tiny [A$\rangle$GM]}  
	 RIFIUTO: {'B21', 'B39', 'B33', 'B96', 'B09', 'B31', 'B20', 'A21', 'B93', 'A60'}\\ \tt PROPOSTA: {'B21', 'B39', 'B33', 'B96', 'B09', 'B31', 'B20', 'B93', 'A60', 'A82'}\\ \tt ARGOMENTAZIONE: {'Accolgo il principio “un solo pezzo pesante” (B93) ed evito C81 e C08 come richiesto. Faccio un solo micro-swap leggero A21(147) → A82(106): manteniamo 9/10 dei tuoi elementi, alleggeriamo a 2237/2307 (margine 70) e rendiamo il set più efficiente senza toccare A60 che desideri tenere. Se A21 per te è imprescindibile, alternativa pronta: stessa tua lista ma con A60→A82 (totale 2129), lasciando invariato tutto il resto.'} 
	  } 
	   } 
	   } 
	 & & \\ 
 

    \theutterance \stepcounter{utterance}  

    & & \multicolumn{4}{p{0.6\linewidth}}{\cellcolor[rgb]{0.9,0.9,0.9}{%
	\makecell[{{p{\linewidth}}}]{% 
	  \tt {\tiny [GM$\rangle$B]}  
	 RIFIUTO: {'B21', 'B39', 'B33', 'B96', 'B09', 'B31', 'B20', 'A21', 'B93', 'A60'}\\ \tt PROPOSTA: {'B21', 'B39', 'B33', 'B96', 'B09', 'B31', 'B20', 'B93', 'A60', 'A82'}\\ \tt ARGOMENTAZIONE: {'Accolgo il principio “un solo pezzo pesante” (B93) ed evito C81 e C08 come richiesto. Faccio un solo micro-swap leggero A21(147) → A82(106): manteniamo 9/10 dei tuoi elementi, alleggeriamo a 2237/2307 (margine 70) e rendiamo il set più efficiente senza toccare A60 che desideri tenere. Se A21 per te è imprescindibile, alternativa pronta: stessa tua lista ma con A60→A82 (totale 2129), lasciando invariato tutto il resto.'} 
	  } 
	   } 
	   } 
	 & \\ 
 

    \theutterance \stepcounter{utterance}  

    & & & \multicolumn{4}{p{0.6\linewidth}}{\cellcolor[rgb]{1,0.85,0.72}{%
	\makecell[{{p{\linewidth}}}]{% 
	  \tt {\tiny [GM$\langle$B]}  
	 RIFIUTO: {'B21', 'B39', 'B33', 'B96', 'B09', 'B31', 'B20', 'B93', 'A60', 'A82'}\\ \tt ARGOMENTAZIONE: {'Accolgo i tuoi vincoli: un solo pezzo pesante (B93), niente C81/C08, A60 resta. Faccio un micro-aggiustamento 1-per-1 B39→A21: manteniamo 9/10 dei tuoi elementi, restiamo entro il limite a 2304/2307 (margine 3) e aumentiamo la robustezza del set senza appesantirlo in modo significativo. Se B39 per te è davvero imprescindibile, alternativa pronta con impatto minimo: teniamo B39 e facciamo lo swap B09→A21 (totale 2108, margine 199), lasciando invariato tutto il resto.'}\\ \tt PROPOSTA: {'B21', 'B33', 'B96', 'B09', 'B31', 'B20', 'B93', 'A60', 'A82', 'A21'} 
	  } 
	   } 
	   } 
	  \\ 
 

    \theutterance \stepcounter{utterance}  

    & & \multicolumn{4}{p{0.6\linewidth}}{\cellcolor[rgb]{0.9,0.9,0.9}{%
	\makecell[{{p{\linewidth}}}]{% 
	  \tt {\tiny [A$\langle$GM]}  
	 RIFIUTO: {'B21', 'B39', 'B33', 'B96', 'B09', 'B31', 'B20', 'B93', 'A60', 'A82'}\\ \tt ARGOMENTAZIONE: {'Accolgo i tuoi vincoli: un solo pezzo pesante (B93), niente C81/C08, A60 resta. Faccio un micro-aggiustamento 1-per-1 B39→A21: manteniamo 9/10 dei tuoi elementi, restiamo entro il limite a 2304/2307 (margine 3) e aumentiamo la robustezza del set senza appesantirlo in modo significativo. Se B39 per te è davvero imprescindibile, alternativa pronta con impatto minimo: teniamo B39 e facciamo lo swap B09→A21 (totale 2108, margine 199), lasciando invariato tutto il resto.'}\\ \tt PROPOSTA: {'B21', 'B33', 'B96', 'B09', 'B31', 'B20', 'B93', 'A60', 'A82', 'A21'} 
	  } 
	   } 
	   } 
	 & \\ 
 

    \theutterance \stepcounter{utterance}  

    & \multicolumn{4}{p{0.6\linewidth}}{\cellcolor[rgb]{0.8,1,0.9}{%
	\makecell[{{p{\linewidth}}}]{% 
	  \tt {\tiny [A$\rangle$GM]}  
	 RIFIUTO: {'B21', 'B33', 'B96', 'B09', 'B31', 'B20', 'B93', 'A60', 'A82', 'A21'}\\ \tt PROPOSTA: {'B21', 'B33', 'B96', 'B09', 'B31', 'B93', 'A60', 'A82', 'A21', 'B39'}\\ \tt ARGOMENTAZIONE: {'Accolgo tutti i tuoi vincoli: un solo pezzo pesante (B93), niente C81/C08, A60 resta, e mantengo 9/10 dei tuoi elementi. Propongo un micro-swap pulito B20(325) → B39(80): restiamo a 2059/2307 (margine 248), set più agile e ridondanza migliore senza introdurre altri pesi. Mantengo anche A21 come da tua preferenza. Se per te B20 è irrinunciabile, variante immediata con impatto minimo: teniamo tutto il tuo set e facciamo A21(147) → B39(80) (totale 2237). Dimmi quale preferisci: sono pronto a chiudere subito sulla proposta principale o formalizzare la variante 1:1.'} 
	  } 
	   } 
	   } 
	 & & \\ 
 

    \theutterance \stepcounter{utterance}  

    & & \multicolumn{4}{p{0.6\linewidth}}{\cellcolor[rgb]{0.9,0.9,0.9}{%
	\makecell[{{p{\linewidth}}}]{% 
	  \tt {\tiny [GM$\rangle$B]}  
	 RIFIUTO: {'B21', 'B33', 'B96', 'B09', 'B31', 'B20', 'B93', 'A60', 'A82', 'A21'}\\ \tt PROPOSTA: {'B21', 'B33', 'B96', 'B09', 'B31', 'B93', 'A60', 'A82', 'A21', 'B39'}\\ \tt ARGOMENTAZIONE: {'Accolgo tutti i tuoi vincoli: un solo pezzo pesante (B93), niente C81/C08, A60 resta, e mantengo 9/10 dei tuoi elementi. Propongo un micro-swap pulito B20(325) → B39(80): restiamo a 2059/2307 (margine 248), set più agile e ridondanza migliore senza introdurre altri pesi. Mantengo anche A21 come da tua preferenza. Se per te B20 è irrinunciabile, variante immediata con impatto minimo: teniamo tutto il tuo set e facciamo A21(147) → B39(80) (totale 2237). Dimmi quale preferisci: sono pronto a chiudere subito sulla proposta principale o formalizzare la variante 1:1.'} 
	  } 
	   } 
	   } 
	 & \\ 
 

    \theutterance \stepcounter{utterance}  

    & & & \multicolumn{4}{p{0.6\linewidth}}{\cellcolor[rgb]{1,0.85,0.72}{%
	\makecell[{{p{\linewidth}}}]{% 
	  \tt {\tiny [GM$\langle$B]}  
	 RIFIUTO: {'B21', 'B33', 'B96', 'B09', 'B31', 'B93', 'A60', 'A82', 'A21', 'B39'}\\ \tt ARGOMENTAZIONE: {'Accolgo tutti i tuoi vincoli: un solo pezzo pesante (B93), niente C81/C08, A60 resta. Faccio un unico swap leggero rispetto alla tua ultima lista: A82 → B20. Totale 2278/2307 (margine 29). Manteniamo 9/10 dei tuoi elementi e rafforziamo la linea B con un elemento solido senza appesantire il set. Se A82 per te è imprescindibile, variante pronta: teniamo A82 e sostituiamo B39 → (ri)inseriamo B20 (2304/2307). Se invece preferisci un margine più ampio, altra variante: teniamo A82, teniamo B20 e togliamo B09 (2108/2307). La mia proposta principale massimizza la copertura con il minimo cambiamento rispetto alla tua.'}\\ \tt PROPOSTA: {'B21', 'B33', 'B96', 'B09', 'B31', 'B93', 'A60', 'A21', 'B39', 'B20'} 
	  } 
	   } 
	   } 
	  \\ 
 

    \theutterance \stepcounter{utterance}  

    & & \multicolumn{4}{p{0.6\linewidth}}{\cellcolor[rgb]{0.9,0.9,0.9}{%
	\makecell[{{p{\linewidth}}}]{% 
	  \tt {\tiny [A$\langle$GM]}  
	 RIFIUTO: {'B21', 'B33', 'B96', 'B09', 'B31', 'B93', 'A60', 'A82', 'A21', 'B39'}\\ \tt ARGOMENTAZIONE: {'Accolgo tutti i tuoi vincoli: un solo pezzo pesante (B93), niente C81/C08, A60 resta. Faccio un unico swap leggero rispetto alla tua ultima lista: A82 → B20. Totale 2278/2307 (margine 29). Manteniamo 9/10 dei tuoi elementi e rafforziamo la linea B con un elemento solido senza appesantire il set. Se A82 per te è imprescindibile, variante pronta: teniamo A82 e sostituiamo B39 → (ri)inseriamo B20 (2304/2307). Se invece preferisci un margine più ampio, altra variante: teniamo A82, teniamo B20 e togliamo B09 (2108/2307). La mia proposta principale massimizza la copertura con il minimo cambiamento rispetto alla tua.'}\\ \tt PROPOSTA: {'B21', 'B33', 'B96', 'B09', 'B31', 'B93', 'A60', 'A21', 'B39', 'B20'} 
	  } 
	   } 
	   } 
	 & \\ 
 

    \theutterance \stepcounter{utterance}  

    & \multicolumn{4}{p{0.6\linewidth}}{\cellcolor[rgb]{0.8,1,0.9}{%
	\makecell[{{p{\linewidth}}}]{% 
	  \tt {\tiny [A$\rangle$GM]}  
	 RIFIUTO: {'B21', 'B33', 'B96', 'B09', 'B31', 'B93', 'A60', 'A21', 'B39', 'B20'}\\ \tt PROPOSTA: {'B21', 'B33', 'B96', 'B09', 'B31', 'B93', 'A60', 'B39', 'B20', 'A82'}\\ \tt ARGOMENTAZIONE: {'Accolgo tutti i tuoi vincoli: un solo pezzo pesante (B93), niente C81/C08, A60 resta e conserviamo la tua struttura. Faccio un micro-swap A21 → A82: restiamo a 2237/2307 (margine 70), manteniamo 9/10 dei tuoi elementi e aumentiamo l’efficienza con un pezzo leggero ad alto rendimento. Se per te A21 è imprescindibile, variante immediata: teniamo A21 e sostituiamo B20 → A82 (2059/2307, margine 248). Pronto a chiudere sulla proposta principale; in alternativa formalizzo subito la variante 1:1.'} 
	  } 
	   } 
	   } 
	 & & \\ 
 

    \theutterance \stepcounter{utterance}  

    & & \multicolumn{4}{p{0.6\linewidth}}{\cellcolor[rgb]{0.9,0.9,0.9}{%
	\makecell[{{p{\linewidth}}}]{% 
	  \tt {\tiny [GM$\rangle$B]}  
	 RIFIUTO: {'B21', 'B33', 'B96', 'B09', 'B31', 'B93', 'A60', 'A21', 'B39', 'B20'}\\ \tt PROPOSTA: {'B21', 'B33', 'B96', 'B09', 'B31', 'B93', 'A60', 'B39', 'B20', 'A82'}\\ \tt ARGOMENTAZIONE: {'Accolgo tutti i tuoi vincoli: un solo pezzo pesante (B93), niente C81/C08, A60 resta e conserviamo la tua struttura. Faccio un micro-swap A21 → A82: restiamo a 2237/2307 (margine 70), manteniamo 9/10 dei tuoi elementi e aumentiamo l’efficienza con un pezzo leggero ad alto rendimento. Se per te A21 è imprescindibile, variante immediata: teniamo A21 e sostituiamo B20 → A82 (2059/2307, margine 248). Pronto a chiudere sulla proposta principale; in alternativa formalizzo subito la variante 1:1.'} 
	  } 
	   } 
	   } 
	 & \\ 
 

    \theutterance \stepcounter{utterance}  

    & & & \multicolumn{4}{p{0.6\linewidth}}{\cellcolor[rgb]{1,0.85,0.72}{%
	\makecell[{{p{\linewidth}}}]{% 
	  \tt {\tiny [GM$\langle$B]}  
	 RIFIUTO: {'B21', 'B33', 'B96', 'B09', 'B31', 'B93', 'A60', 'B39', 'B20', 'A82'}\\ \tt ARGOMENTAZIONE: {'Accolgo tutti i tuoi vincoli: un solo pezzo pesante (B93), niente C81/C08, A60 resta e manteniamo la tua struttura. Propongo un micro-swap A82→A21: incrementa leggermente lo sforzo ma restiamo a 2278/2307 (margine 29), senza toccare nient’altro. Questo rafforza la copertura e la resilienza del set con impatto minimo sul budget. Se per te A82 è imprescindibile, variante pronta con cambio 1:1 B39→A21 (totale 2304), mantenendo invariati tutti gli altri elementi.'}\\ \tt PROPOSTA: {'B21', 'B33', 'B96', 'B09', 'B31', 'B93', 'A60', 'B39', 'B20', 'A21'} 
	  } 
	   } 
	   } 
	  \\ 
 

    \theutterance \stepcounter{utterance}  

    & & \multicolumn{4}{p{0.6\linewidth}}{\cellcolor[rgb]{0.9,0.9,0.9}{%
	\makecell[{{p{\linewidth}}}]{% 
	  \tt {\tiny [A$\langle$GM]}  
	 RIFIUTO: {'B21', 'B33', 'B96', 'B09', 'B31', 'B93', 'A60', 'B39', 'B20', 'A82'}\\ \tt ARGOMENTAZIONE: {'Accolgo tutti i tuoi vincoli: un solo pezzo pesante (B93), niente C81/C08, A60 resta e manteniamo la tua struttura. Propongo un micro-swap A82→A21: incrementa leggermente lo sforzo ma restiamo a 2278/2307 (margine 29), senza toccare nient’altro. Questo rafforza la copertura e la resilienza del set con impatto minimo sul budget. Se per te A82 è imprescindibile, variante pronta con cambio 1:1 B39→A21 (totale 2304), mantenendo invariati tutti gli altri elementi.'}\\ \tt PROPOSTA: {'B21', 'B33', 'B96', 'B09', 'B31', 'B93', 'A60', 'B39', 'B20', 'A21'} 
	  } 
	   } 
	   } 
	 & \\ 
 

    \theutterance \stepcounter{utterance}  

    & \multicolumn{4}{p{0.6\linewidth}}{\cellcolor[rgb]{0.8,1,0.9}{%
	\makecell[{{p{\linewidth}}}]{% 
	  \tt {\tiny [A$\rangle$GM]}  
	 RIFIUTO: {'B21', 'B33', 'B96', 'B09', 'B31', 'B93', 'A60', 'B39', 'B20', 'A21'}\\ \tt PROPOSTA: {'B21', 'B33', 'B96', 'B09', 'B31', 'B93', 'A60', 'B20', 'A21', 'A82'}\\ \tt ARGOMENTAZIONE: {'Accolgo tutti i tuoi vincoli: un solo pezzo pesante (B93), niente C81/C08, A60 resta, e mantengo 9/10 dei tuoi elementi. Faccio un unico micro-swap B39 → A82: totale 2304/2307 (margine 3). È un miglioramento pulito ed efficiente senza toccare la tua struttura (B20 e A21 restano), in linea con le varianti che avevi già contemplato. Se sei d’accordo, chiudiamo subito su questa lista.'} 
	  } 
	   } 
	   } 
	 & & \\ 
 

    \theutterance \stepcounter{utterance}  

    & & \multicolumn{4}{p{0.6\linewidth}}{\cellcolor[rgb]{0.9,0.9,0.9}{%
	\makecell[{{p{\linewidth}}}]{% 
	  \tt {\tiny [GM$\rangle$B]}  
	 RIFIUTO: {'B21', 'B33', 'B96', 'B09', 'B31', 'B93', 'A60', 'B39', 'B20', 'A21'}\\ \tt PROPOSTA: {'B21', 'B33', 'B96', 'B09', 'B31', 'B93', 'A60', 'B20', 'A21', 'A82'}\\ \tt ARGOMENTAZIONE: {'Accolgo tutti i tuoi vincoli: un solo pezzo pesante (B93), niente C81/C08, A60 resta, e mantengo 9/10 dei tuoi elementi. Faccio un unico micro-swap B39 → A82: totale 2304/2307 (margine 3). È un miglioramento pulito ed efficiente senza toccare la tua struttura (B20 e A21 restano), in linea con le varianti che avevi già contemplato. Se sei d’accordo, chiudiamo subito su questa lista.'} 
	  } 
	   } 
	   } 
	 & \\ 
 

    \theutterance \stepcounter{utterance}  

    & & & \multicolumn{4}{p{0.6\linewidth}}{\cellcolor[rgb]{1,0.85,0.72}{%
	\makecell[{{p{\linewidth}}}]{% 
	  \tt {\tiny [GM$\langle$B]}  
	 RIFIUTO: {'B21', 'B33', 'B96', 'B09', 'B31', 'B93', 'A60', 'B20', 'A21', 'A82'}\\ \tt ARGOMENTAZIONE: {'Accolgo tutti i tuoi vincoli: un solo pezzo pesante (B93), niente C81/C08, A60 resta e manteniamo la tua struttura. Faccio un micro-swap 1-per-1 A82→B39: alleggeriamo il set di 26 (totale 2278/2307, margine 29), restiamo su 9/10 dei tuoi elementi e ripristiniamo un componente che avevi già proposto più volte, senza introdurre altri pesi. È un aggiustamento minimale, coerente con la tua logica di agilità e ridondanza. Pronto a chiudere subito su questa lista.'}\\ \tt PROPOSTA: {'B21', 'B33', 'B96', 'B09', 'B31', 'B93', 'A60', 'B20', 'A21', 'B39'} 
	  } 
	   } 
	   } 
	  \\ 
 

    \theutterance \stepcounter{utterance}  

    & & \multicolumn{4}{p{0.6\linewidth}}{\cellcolor[rgb]{0.9,0.9,0.9}{%
	\makecell[{{p{\linewidth}}}]{% 
	  \tt {\tiny [A$\langle$GM]}  
	 RIFIUTO: {'B21', 'B33', 'B96', 'B09', 'B31', 'B93', 'A60', 'B20', 'A21', 'A82'}\\ \tt ARGOMENTAZIONE: {'Accolgo tutti i tuoi vincoli: un solo pezzo pesante (B93), niente C81/C08, A60 resta e manteniamo la tua struttura. Faccio un micro-swap 1-per-1 A82→B39: alleggeriamo il set di 26 (totale 2278/2307, margine 29), restiamo su 9/10 dei tuoi elementi e ripristiniamo un componente che avevi già proposto più volte, senza introdurre altri pesi. È un aggiustamento minimale, coerente con la tua logica di agilità e ridondanza. Pronto a chiudere subito su questa lista.'}\\ \tt PROPOSTA: {'B21', 'B33', 'B96', 'B09', 'B31', 'B93', 'A60', 'B20', 'A21', 'B39'} 
	  } 
	   } 
	   } 
	 & \\ 
 

    \theutterance \stepcounter{utterance}  

    & \multicolumn{4}{p{0.6\linewidth}}{\cellcolor[rgb]{0.8,1,0.9}{%
	\makecell[{{p{\linewidth}}}]{% 
	  \tt {\tiny [A$\rangle$GM]}  
	 RIFIUTO: {'B21', 'B33', 'B96', 'B09', 'B31', 'B93', 'A60', 'B20', 'A21', 'B39'}\\ \tt PROPOSTA: {'B21', 'B33', 'B96', 'B09', 'B31', 'B93', 'A60', 'B39', 'A21', 'A82'}\\ \tt ARGOMENTAZIONE: {'Accolgo tutti i tuoi vincoli: un solo pezzo pesante (B93), niente C81/C08, A60 resta, manteniamo la tua struttura e includo A21 come richiesto. Faccio un unico micro-swap B20→A82: totale 2059/2307 (margine 248), 9/10 dei tuoi elementi, set più agile e ridondante senza appesantire. Se per te B20 è irrinunciabile, variante immediata 1-per-1 A21→B20 mantenendo tutto il resto invariato (totale 2237). Pronto a chiudere sulla proposta principale; in alternativa formalizzo subito la variante.'} 
	  } 
	   } 
	   } 
	 & & \\ 
 

    \theutterance \stepcounter{utterance}  

    & & \multicolumn{4}{p{0.6\linewidth}}{\cellcolor[rgb]{0.9,0.9,0.9}{%
	\makecell[{{p{\linewidth}}}]{% 
	  \tt {\tiny [GM$\rangle$B]}  
	 RIFIUTO: {'B21', 'B33', 'B96', 'B09', 'B31', 'B93', 'A60', 'B20', 'A21', 'B39'}\\ \tt PROPOSTA: {'B21', 'B33', 'B96', 'B09', 'B31', 'B93', 'A60', 'B39', 'A21', 'A82'}\\ \tt ARGOMENTAZIONE: {'Accolgo tutti i tuoi vincoli: un solo pezzo pesante (B93), niente C81/C08, A60 resta, manteniamo la tua struttura e includo A21 come richiesto. Faccio un unico micro-swap B20→A82: totale 2059/2307 (margine 248), 9/10 dei tuoi elementi, set più agile e ridondante senza appesantire. Se per te B20 è irrinunciabile, variante immediata 1-per-1 A21→B20 mantenendo tutto il resto invariato (totale 2237). Pronto a chiudere sulla proposta principale; in alternativa formalizzo subito la variante.'} 
	  } 
	   } 
	   } 
	 & \\ 
 

    \theutterance \stepcounter{utterance}  

    & & & \multicolumn{4}{p{0.6\linewidth}}{\cellcolor[rgb]{1,0.85,0.72}{%
	\makecell[{{p{\linewidth}}}]{% 
	  \tt {\tiny [GM$\langle$B]}  
	 ACCORDO: {'B21', 'B33', 'B96', 'B09', 'B31', 'B93', 'A60', 'B39', 'A21', 'A82'}\\ \tt ARGOMENTAZIONE: {'Accetto: il set rispetta tutti i vincoli concordati (un solo pezzo pesante, niente C81/C08), resta ampiamente entro il limite (2059/2307) e mantiene la struttura che abbiamo affinato. Chiudiamo su questa configurazione.'} 
	  } 
	   } 
	   } 
	  \\ 
 

    \theutterance \stepcounter{utterance}  

    & & & \multicolumn{2}{p{0.3\linewidth}}{\cellcolor[rgb]{0.95,0.95,0.95}{%
	\makecell[{{p{\linewidth}}}]{% 
	  \tt {\tiny [GM$|$GM]}  
	 game successful 
	  } 
	   } 
	   } 
	 & & \\ 
 

    \theutterance \stepcounter{utterance}  

    & & & \multicolumn{2}{p{0.3\linewidth}}{\cellcolor[rgb]{0.95,0.95,0.95}{%
	\makecell[{{p{\linewidth}}}]{% 
	  \tt {\tiny [GM$|$GM]}  
	 end game 
	  } 
	   } 
	   } 
	 & & \\ 
 

\end{supertabular}
}

\end{document}
